The term follows the limited-attention literature, in which an agent who does not attend to an attribute replaces it with a default value rather than with an inference \citep{gabaix_sparsity_2014, gabaix_behavioral_2019}. The default is therefore not a belief about the dwelling's rating. It is what the buyer implicitly assumes when the question of energy never arises. Three benchmark values fix ideas. With $C^{\rm def}=0$ no energy cost is priced at all, which does not mean the buyer believes the dwelling consumes nothing, only that energy plays no part in her valuation. With $C^{\rm def}$ equal to the prior mean cost $\bar\theta C_1+(1-\bar\theta)C_0$, the inattentive buyer behaves like one who assumes an average dwelling, which amounts to cursedness. With $C^{\rm def}=\bar C^{\rm sil}$ she happens to subtract exactly what an attentive buyer would infer, and attention becomes irrelevant. The same default applies to rated and silent listings, because a buyer who does not attend to energy, by definition, does not use what the listing shows about it. Which value is relevant is an empirical question: Proposition~\ref{prop:att_incentives}(iii) shows that the extension can explain the apartment wedge only if $C^{\rm def}<C_1$, and Proposition~\ref{prop:att_identification}(iii) shows what the data can say about its level.
